\begin{table}[H]
\centering
\caption{\label{tab:open-ended-categories}Category-specific reliability of primary open-ended codes}
\centering
\fontsize{9}{11}\selectfont
\begin{threeparttable}
\begin{tabular}[t]{rlrrr}
\toprule
Code & Category & Anne n & Jeremy n & Binary kappa\\
\midrule
1 & No stated reason & 169 & 170 & 0.882\\
2 & Equal expected value & 162 & 142 & 0.824\\
3 & Certainty or risk aversion & 582 & 636 & 0.786\\
4 & Risk offers chance all survive & 726 & 718 & 0.760\\
5 & Saves more people & 443 & 349 & 0.638\\
6 & Easier to understand & 12 & 6 & 0.331\\
7 & Other & 137 & 163 & 0.508\\
8 & Better choice without a reason & 64 & 111 & 0.603\\
\bottomrule
\end{tabular}
\begin{tablenotes}[para]
\item Note: Counts and binary kappa values use primary codes only. The instructions allowed multiple categories. A contemporaneous memo identifies overlap among codes 3, 4, and 5 and ambiguity involving codes 1, 7, and 8.
\end{tablenotes}
\end{threeparttable}
\end{table}
